\section{Learning to wait in the company of the saints}

\subsection{Who is coming to judge us?}

Imagine standing with the people in Jerusalem who heard St Cyril explain
the Creed in the fourth century. The judge of the living and the dead is
the man who once stood before a judge and scarcely opened his mouth. How
will we know him? Cyril puts two garments beside one another. A mother
wrapped her son in strips of cloth. At his coming in glory, light itself
will be his garment. He who went to the Cross in our flesh will return with
the angels. We have already been shown his patience. We still await the
day when his kingship will be plain to every eye.

One could spend a lifetime wondering when that day will come. Cyril sends
his hearers to Matthew 25 instead. The Lord who will pronounce judgment
tells us where he can be served today. At the end of the same instruction
Cyril gathers two Gospel pictures together, the hungry guest and the
bridesmaids waiting with their lamps.

\begin{quote}
At once then begin to do these works, and abide in the faith; lest, like
the foolish virgins, tarrying to buy oil, thou be shut out. Be not confident
because thou merely possessest the lamp, but constantly keep it burning.
\end{quote}

Who is waiting for a visit? A person who has grown frail can keep watch too.
A telephone call to another sick person can be an act of mercy in both
directions. Cyril's lecture was instruction for people preparing for baptism.
Its first paragraphs have since become a reading in the Advent Office. The
instruction about the lamp is farther on in the lecture.\footnote{Cyril of
Jerusalem, \emph{Catechetical Lecture} 15.1--4, 26, trans. E.~H. Gifford,
\emph{Nicene and Post-Nicene Fathers}, second series, VII (1894), pp.~104--105,
112, \href{https://www.ccel.org/ccel/schaff/npnf207.html}{English edition}.
\emph{Liturgy of the Hours}, vol.~I, First Sunday of Advent, Office of
Readings, second reading, from Lecture 15.1--3.}

\subsection{What becomes of a word we hear at Mass?}

St Bernard of Clairvaux preached to monks who sang the Scriptures for
hours each day. They needed to hear Christ's promise in John 14:23 as
much as anyone else. The Father and Son will make their home with the
person who loves Christ and keeps his word. Bernard called this the hidden
coming of Christ between the two comings the Church professes in the
Creed. Jesus lived among people who could see and touch him, and he will
come again for all to see. In the meantime a man can stand beside his
brother in choir and be quite unaware of what Christ is doing in that
brother's soul.

How does a person keep a word? Bernard took his monks to their cupboard.
Bread left on the shelf can go stale or be eaten by mice. Bread that has
become part of your body is kept in quite another way.

\begin{quote}
Receive this spiritual bread, therefore, into the stomach of your minds,
and having been there properly digested, let it be transmitted thence to
the affections and the will.
\end{quote}

Go back once today to a sentence you heard at Mass. If it was John's call
to share food, stay with it long enough for someone you know to come to
mind. Desire that person's good. Give what the person can actually use.
There is Bernard's movement from hearing to affection to the will. He
promises something more in the last part of his sermon. A word kept in
this way begins to keep the hearer. The Son comes with his Father to
rebuild the holy city within. The Roman Office reads from this sermon on
Wednesday of the first week.\footnote{Bernard of Clairvaux, Fifth Sermon for
Advent, nos.~1--3, trans. a Priest of Mount Melleray,
\emph{St Bernard's Sermons for the Seasons}, I (Dublin, 1921),
\href{https://archive.org/details/stbernardssermon0001prie_y4z3/page/40/mode/2up}{pp.~40--43}.
Congregation for Divine Worship and the Discipline of the Sacraments,
\href{https://www.usccb.org/resources/Homiletic-Directory-CDWDS.pdf}{\emph{Homiletic Directory}},
no.~79.}

Think of the chalice at Mass. People who have come to honor Christ's Body
want the vessel to be worthy of what it holds. St John Chrysostom, in a
homily on Matthew, takes that desire seriously enough to ask where the money
for such a gift came from. A gift bought by cheating a widow or an orphan
brings that injustice right to the altar. He next takes his congregation
into the street. The speaker of the words of institution is the very person
who says that the hungry, the naked and the prisoner have a claim on us.

\begin{quote}
Dost thou make Him a cup of gold, while thou givest Him not a cup of cold water?
\end{quote}

Chrysostom approves of giving beautiful things for divine worship from honest
earnings. He insists on feeding the hungry first. The one who drinks the water
receives a real good, as does the person who gives it. When the parish sets
out its collection for families in need, buy the items the families have
asked for. Here is a way of answering the King in Matthew 25 before we find
ourselves face to face with him.\footnote{John Chrysostom,
\emph{Homilies on Matthew}, 50.4--5, on Matthew 14, drawing also on
Matthew 25, trans. George Prevost, rev. M.~B. Riddle, \emph{Nicene and
Post-Nicene Fathers}, first series, X (1888),
\href{https://www.ccel.org/ccel/schaff/npnf110.html}{English edition}.}

\subsection{When the bucket comes up empty}

Proba, a widow with a family, asked St Augustine how she should pray. In
his letter he took up a difficulty any Christian who has repeatedly prayed
for someone will recognize. The Father already knows what we need. Why
return to the same petition? Augustine looked at the distance between what
God wishes to give and what a human heart can take in. In his words, “His
gifts are very great, but we are small and straitened in our capacity of
receiving.” Asking exercises the desire by which we receive. In the course
of a day or a long life that desire can become larger. Faith, hope and
charity can also persist through hours spent doing something else. Returning
to the actual words of prayer at a set time reawakens the desire that
ordinary cares have cooled. Try it at the end of the day with one petition
of the Our Father. Consider whose bread you have been asking the Father
to give, and whom you have included in the word \emph{our}.
\footnote{Augustine, Letter 130 to Proba, nos.~17--21, trans.
J.~G. Cunningham, \emph{Nicene and Post-Nicene Fathers}, first series, I
(1886), \href{https://www.ccel.org/ccel/schaff/npnf101.html}{pp.~465--467}.}

St Teresa of Ávila describes the effort of such prayer from the inside.
In her \emph{Life} she thinks of the soul as a garden entrusted to someone
who has to draw water by hand. Lower the bucket. Haul it up. Sometimes it
comes up empty. The reader who has gone to prayer for days without being
able to hold a thought or shed a tear will know the weight of that bucket.
Teresa calls tears or tenderness the water. The flowers are the virtues
God grows. We can give the Lord the time we promised, desire his good
pleasure and accompany him in carrying the Cross. The water in the well
belongs to him.

\begin{quote}
God grant there may be water in it! That, however, does not depend on us;
we are drawing it, and doing what we can towards watering the flowers.
\end{quote}

A family can sing the Advent hymn in the kitchen and someone at the table
can feel only grief. Let that person sing or keep silence as strength
allows. Teresa believes that even in the dry garden the Lord can preserve
what is growing. A person can still give him the care of what he has
entrusted to them today.\footnote{Teresa of Ávila, \emph{Life} 11.10--16,
trans. David Lewis, third edition with notes by Benedict Zimmerman (London,
1904), \href{https://archive.org/details/lifeofteresajesu00tereuoft/page/78/mode/2up}{pp.~78--81},
especially p.~80.}

\subsection{Why the Church waits with Mary}

St Paul VI invites a congregation to listen twice to Luke's account of
the Visitation. First Elizabeth speaks. She blesses the woman who believed
that the Lord's word would be fulfilled (Luke 1:45). Then Mary speaks.
Her song remembers what God has done for Israel and what he promised to
Abraham. When the Church sings the Magnificat in the evening, she receives
both a prayer and a way of praying. She hears the saving deed before she
tries to speak of it. Then she borrows the words of the woman who carried
the Savior to another woman's home. Paul VI goes on to Cana, where Mary
takes a household's quite ordinary embarrassment to her Son, and to the
upper room, where she prays with the apostles. A Christian keeping the last
week of Advent can do likewise. Listen to the Gospel of the day. Praise
God for a kindness received. Tell Christ about a need in a house other
than your own. Ask his mother to pray with you.

Paul VI made this argument in an apostolic exhortation about devotion to
Mary. He singled out Advent, especially the days from December 17 and the
Sunday before Christmas. In the last days before Christmas, Mary teaches the Church how to ask her
Son to come. She already knows what it means to wait for him.
\footnote{Paul VI,
\href{https://www.vatican.va/content/paul-vi/en/apost_exhortations/documents/hf_p-vi_exh_19740202_marialis-cultus.html}{\emph{Marialis cultus}},
February 2, 1974, nos.~3--4, 17--18. The account here is a paraphrase
of the exhortation.}

\subsection{At home after Midnight Mass}

On Christmas night in 1886, a girl of thirteen returned from Midnight Mass
to a familiar family custom. Her shoes stood by the hearth filled with
gifts. Her father had always liked to see her open them. This time Thérèse
overheard him say that he hoped this would be the last year of it. Her
sister Céline thought she would cry and urged her to wait upstairs. Thérèse
felt the hurt. She also found she could go down immediately, pick up her
shoes and share the delight of opening the presents with her father. He
laughed. Céline could hardly believe what she had seen.

Years later the Carmelite nun told the story to her sister in religion.
She called it her conversion. In the Child who had accepted human weakness
she recognized the one who had given her the strength to love. Her account
continues with Peter and his empty nets in Luke 5. She had labored over
her faults for years. On that night Christ himself filled the net. The
fruit she remembered was an ability to forget herself for others. Consider
what such a grace could mean for one conversation this Christmas. There
may be someone whose joy you are free to share, even after an ill-chosen
remark. Thérèse's tale gives this hope a place as ordinary as a flight of
stairs.\footnote{Thérèse of Lisieux, \emph{Manuscript A}, 44v--45v,
Christmas 1886. See the \href{https://archives.carmeldelisieux.fr/naissance-dune-sainte/iconographie-theresienne/la-vie-en-images/du-sourire-de-la-vierge-a-lentree-au-carmel/}{Carmel of Lisieux},
and the historical English version, \emph{Sœur Thérèse of Lisieux}, ed.
T.~N. Taylor (London, Burns \& Oates, fifteenth thousand, imprimatur
1912), chap.~V, \href{https://archive.org/details/saintthaeraeseof00thaeuoft/page/71/mode/2up}{pp.~71--72}.
The episode is retold here in paraphrase.}

Long before that French household kept Christmas, St Leo spoke in Rome
to three kinds of hearers standing before the same mystery. A saint could
rejoice because victory was near. A sinner could rejoice because pardon
was offered. A person who had yet to believe could take heart because
Christ was calling him to life. The reason for their joy was the same for
all three. The Son had taken a human nature in order to reconcile it to its
Creator. Leo then addressed the baptized. The body that can rise from a chair
to go and make peace with someone is a member of Christ. Baptism has made
it a dwelling of the Holy Spirit. The debt Christ paid for that person
cost him his blood. That is how much hangs on the tenderness we see in the
crib, and how much a person matters to the Child we have come to adore.
\footnote{Leo the Great, Nativity Sermon 21.1--3, trans. C.~L. Feltoe,
\emph{Nicene and Post-Nicene Fathers}, second series, XII (1895),
\href{https://www.ccel.org/ccel/schaff/npnf212.html}{pp.~128--129}.}
